\documentclass[12pt]{article}
\usepackage[utf8]{inputenc}
\usepackage[T1]{fontenc}
\usepackage[margin=1in]{geometry}
\usepackage{setspace}
\usepackage{xcolor}
\usepackage[colorlinks=true, urlcolor=blue!60!black, linkcolor=blue!60!black]{hyperref}
\setlength{\parindent}{0em}
\setlength{\parskip}{0.6em}
\renewcommand{\baselinestretch}{1.0}
\begin{document}
\begin{center}
{\LARGE \textit{California v. Texas}, 593 U.S. 659 (2021)\par}
\vspace{0.5em}
{\large Case Brief --- Two-Page Edition\par}
\end{center}

\section{Memory Jogger}
The Court dismissed the third major challenge to the Affordable Care Act, 7--2, holding that no plaintiff had standing to attack an individual mandate whose penalty had been reduced to \$0 (\href{https://www.courtlistener.com/opinion/4892571/california-v-texas/}{593 U.S.\ 659}).

\section{Detailed Case Facts}
\textit{NFIB v.\ Sebelius} (2012) upheld the ACA's individual mandate, 26 U.S.C.\ \S\ 5000A(a), as a tax. In 2017 Congress set the penalty at \$0 but left the coverage language in place. Texas, 17 other States, and two individuals sued, arguing that a revenue-free mandate exceeded the Taxing Power and was inseverable from the whole ACA. When the federal government declined to defend, California and other States intervened.

\section{Procedural History}
The district court held the mandate unconstitutional and inseverable from the entire Act. The Fifth Circuit affirmed on the merits but remanded on severability. The Supreme Court granted cross-petitions.

\section{Judicial Votes}
Dismissed, 7--2. Breyer wrote for Roberts, Thomas, Sotomayor, Kagan, Kavanaugh, and Barrett; Thomas also concurred; Alito, joined by Gorsuch, dissented.

\section{Holding}
No plaintiff showed an injury fairly traceable to actual or threatened enforcement of \S\ 5000A(a). Unenforceable statutory text with no penalty cannot by itself support traceability or redressability, so the Court did not reach the merits or severability.

\section{Analysis of Opinions}
Breyer reasoned that with a \$0 penalty no officer could enforce the mandate, so the individuals' premium costs were not traceable to the government; remedies run against officials, not abstract rules, and a declaration would be advisory. The States' claim that the mandate drove enrollment in Medicaid and CHIP rested on unproven third-party choices, and their reporting costs flowed from separate, unchallenged provisions. Standing-through-inseverability was forfeited. Thomas suggested that theory might work in a future case. Alito, dissenting, called the case the third chapter of an ACA trilogy, found concrete state costs, and would have held the mandate unconstitutional and inseverable.

\section{Examples: Future Applications}
\textbf{Same-Side:} a homeowner who buys compost bins to obey a \$0-penalty federal directive lacks standing. \textbf{Opposite-Side:} a \$200 IRS assessment collected by refund offset supplies enforcement a court can enjoin. \textbf{Fence-Sitter:} a utility pleads standing-through-inseverability from the start, relying on an express inseverability clause.

\section{Critique}
The \textit{Harvard Law Review} noted that no opinion addressed \textit{Massachusetts v.\ EPA}'s ``special solicitude'' for States. Progressives welcomed the result but warned of tighter standing; some textualists praised the refusal to issue advisory declarations, while others saw selective rules shielding the ACA.

\section{Key Quotations}
\begin{quote}
``To find standing here to attack an unenforceable statutory provision would allow a federal court to issue what would amount to `an advisory opinion without the possibility of any judicial relief.' ''
\end{quote}
\noindent --- Justice Breyer, for the Court.

\end{document}
